\section*{Chapter \ref{ch:m2:fx-swaps-forwards-and-the-cross-currency-basis} --- FX Swaps, Forwards and the Cross-Currency Basis}

\begin{solution}{exo:m2:fx-swaps-forwards-and-the-cross-currency-basis:1}
$1.1464 \times (1 + 0.0368 \times 90/360)/(1 + 0.02 \times 90/360) = 1.15119$:
$+47.9$ pips. The euro, with the lower rate, trades at a forward premium.
\end{solution}

\begin{solution}{exo:m2:fx-swaps-forwards-and-the-cross-currency-basis:2}
Otherwise there would be a riskless profit: borrow the low-rate currency, convert
it spot, lend at the high rate and sell the proceeds forward. For this to earn
nothing, the forward must give back less of the low-rate currency per unit of
the high-rate one than spot did: the high-rate currency's forward discount offsets
its interest advantage.
\end{solution}

\begin{solution}{exo:m2:fx-swaps-forwards-and-the-cross-currency-basis:3}
The investor has borrowed USD~10 million and lent JPY~1\,568.70 million for three
months, each loan the collateral of the other. It receives back JPY~1\,557.06
million: the JPY~11.64 million difference is the cost of the dollars, the dollar
interest less the yen interest, plus the basis.
\end{solution}

\begin{solution}{exo:m2:fx-swaps-forwards-and-the-cross-currency-basis:4}
$-25$ basis points: lending yen and borrowing dollars through the swap earns the
yen rate less 0.25\%. Dollars are dearer through the swap than in cash, the usual
sign since 2007.
\end{solution}

\begin{solution}{exo:m2:fx-swaps-forwards-and-the-cross-currency-basis:5}
Long dollars, the trader sells them for the near date and buys them back for the
next at 1.18 pips (0.0118 yen) lower: it earns about JPY~1.18 million a day,
the carry of holding the higher-yielding currency.
\end{solution}

\begin{solution}{exo:m2:fx-swaps-forwards-and-the-cross-currency-basis:6}
At quarter ends banks shrink their balance sheets for their reports and
regulatory ratios, so arbitrage capacity falls and those needing dollars through
swaps pay more: the basis widens. In March 2020 swap lines let central banks
lend dollars to their banks at the OIS rate plus 25 basis points; no bank with
access needed to pay much more than that in swaps, so the basis narrowed.
\end{solution}

\begin{solution}{exo:m2:fx-swaps-forwards-and-the-cross-currency-basis:7}
December 2018; 57 of the 101 months in the data.
\end{solution}

\begin{solution}{exo:m2:fx-swaps-forwards-and-the-cross-currency-basis:8}
The trade uses balance sheet: the bank lends dollars and holds yen assets, which
count in its leverage and capital ratios, and at quarter ends cost more still.
It will do the trade only if the basis pays for that; the basis is the price of
balance sheet, not a free profit, and hedging demand keeps it open.
\end{solution}

\begin{solution}{pb:m2:fx-swaps-forwards-and-the-cross-currency-basis:1}
\textbf{1.} 155.8028, $-106.7$ pips.
\textbf{2.} $3.68 - 0.977 = 2.70$ percentage points a year.
\textbf{3.} $4.75 - 2.70 = 2.05\%$.
\textbf{4.} 1.80\% and 1.55\%.
\textbf{5.} All below the JGB's 2.94\%, by 0.89, 1.14 and 1.39 points.
\textbf{6.} It closes the maturing swap, paying or receiving yen for the change in
the dollar's value since the last roll, and opens a new three-month swap at the
new spot and points.
\textbf{7.} The swap gains in yen what the bonds lose: the insurer receives about
10\% of the hedged notional in yen at the roll, offsetting the fall in the yen
value of its dollar bonds.
\textbf{8.} When the dollar rises, the hedge loses and the insurer must pay yen at
the roll while the gain on the bonds is unrealised; it needs cash for that.
\textbf{9.} It fixes the cost for a year and avoids three rolls, but at the one-year
points and basis, and settles the whole year's revaluation at once.
\textbf{10.} Lower dollar rates, higher yen rates, or a narrower basis.
\textbf{11.} 4.75\% plus or minus the change in the dollar, which in a bad year can
be many times the 1.81-point advantage.
\textbf{12.} The dollar must not fall by more than about 1.81\% over the year against
the yen for the unhedged Treasury to match the JGB.
\textbf{13.} To trade some yield against some currency risk, within its limits, and
to reduce its need for rolling cash.
\textbf{14.} They are one-way demand to borrow dollars through swaps; the banks that
supply them charge for balance sheet.
\textbf{15.} Banks, and through them investors with dollars to lend, such as money
funds lending in repo, and at times central-bank swap lines.
\textbf{16.} In October 2022, just before the Federal Reserve's rates had risen far
enough above Japan's to push the zero-basis hedged yield below the JGB's.
\textbf{17.} It lowers the dollar rate and the hedge cost by 1 point: the hedged
yield rises to about 3.05\% with no basis, above the JGB's 2.94\%, if bond yields
do not move.
\textbf{18.} Because it uses the bank's balance sheet, which is costly; the basis
pays for it, and at quarter ends the bank may not lend at all.
\textbf{19.} Named result: \emph{the hedged yield} of the ten-year Treasury for the
insurer was about 2.05\% with no basis, 1.80\% with a $-25$ basis-point basis,
against 2.94\% on the JGB.
\textbf{20.} The foreign yield less the short-rate differential and the basis:
the hedge takes back what the higher rates give.
\end{solution}

\begin{solution}{iq:m2:fx-swaps-forwards-and-the-cross-currency-basis:1}
Borrow one unit of the base currency at $r_b$, sell it spot for $S$, lend the
proceeds at $r_q$, and agree today to buy back the base at $F$: no initial cost,
no risk, so $F(1 + r_b\tau_b) = S(1 + r_q\tau_q)$, that is
$F = S(1 + r_q\tau_q)/(1 + r_b\tau_b)$.
\iqlookfor{the replication argument and the day counts.}
\end{solution}

\begin{solution}{iq:m2:fx-swaps-forwards-and-the-cross-currency-basis:2}
A spot exchange and its reversal at a forward date: a collateralised loan of one
currency against another. Banks fund foreign-currency assets with it, investors
roll currency hedges with it, dealers manage positions and liquidity with it, and
most of it is very short-dated, so it is traded again and again.
\iqlookfor{the loan interpretation and the uses.}
\end{solution}

\begin{solution}{iq:m2:fx-swaps-forwards-and-the-cross-currency-basis:3}
First, in 2008 and in the euro crisis, banks' credit and funding strains stopped
arbitrage. Since 2014, strong one-way hedging demand to borrow dollars through
swaps (investors, issuers swapping foreign bonds, banks) has met limits to
arbitrage: regulation makes balance sheet costly, especially at quarter ends,
so the gap persists at the price of balance sheet.
\iqlookfor{demand plus limits to arbitrage, not a single cause.}
\end{solution}

\begin{solution}{iq:m2:fx-swaps-forwards-and-the-cross-currency-basis:4}
Sell euros forward against dollars for the portfolio's value, rolling one- or
three-month \omterm{def:m2:fx-swaps-forwards-and-the-cross-currency-basis:swap}{FX swaps} and adjusting the notional as the value changes. The hedge
earns the dollar rate less the euro rate, plus or minus the basis: with dollar
rates higher, a dollar investor gains from hedging euro bonds, and the negative
euro basis adds to that gain. Rebalancing and cash flows at each roll are the
operational costs.
\iqlookfor{the sign of the carry and the role of the basis.}
\end{solution}

\begin{solution}{iq:m2:fx-swaps-forwards-and-the-cross-currency-basis:5}
Swaps that span the quarter end become dearer: banks that report on that date
shrink the balance sheet they lend, so implied dollar rates and the basis jump.
A bank with spare balance sheet lends into the spike; one that needs dollars
funds ahead of time, before the turn is priced.
\iqlookfor{the balance-sheet mechanism and a pre-funding response.}
\end{solution}

\begin{solution}{iq:m2:fx-swaps-forwards-and-the-cross-currency-basis:6}
For each currency, a discount or money-market curve refreshed from its sources;
for each pair, spot and market points by tenor. Compute CIP points from the curves
with the right day counts and spot dates, and the implied basis from market
points; interpolate broken dates on the basis, not on the points; flag stale
inputs and outliers; publish with timestamps; test against hand-computed cases and
round trips between points and basis.
\iqlookfor{conventions per currency, interpolation on the basis, and checks.}
\end{solution}
